% Hebrew research summary for topo-pilot. Build: python build.py (LuaLaTeX, two passes).
\documentclass[12pt,a4paper]{article}

% ---- language, fonts ----------------------------------------------------------------------
\usepackage[bidi=basic,layout=tabular.lists.graphics]{babel}
\babelprovide[main,import]{hebrew}
\babelprovide[import]{english}
\babelfont{rm}[Path=./fonts/,Extension=.ttf,UprightFont=*-Regular,BoldFont=*-Bold]{DavidLibre}
\babelfont{tt}[Scale=0.85]{Consolas}
\newcommand{\en}[1]{\foreignlanguage{english}{#1}}   % isolate an English phrase
\DeclareRobustCommand{\code}[1]{\en{\texttt{\detokenize{#1}}}}

% ---- page ---------------------------------------------------------------------------------
\usepackage[a4paper,top=2.5cm,bottom=2.3cm,left=2.2cm,right=2.2cm,headsep=0.8cm]{geometry}
\usepackage{amsmath,amssymb}
\usepackage{graphicx}
\usepackage{booktabs,tabularx,array}
\usepackage{xcolor}
\usepackage[most]{tcolorbox}
\usepackage{enumitem}
\usepackage[font=small,labelfont=bf,justification=justified,singlelinecheck=false]{caption}
\usepackage[section]{placeins}
\usepackage{float}
\makeatletter\setlength{\@fptop}{0pt}\setlength{\@fpsep}{12pt}\makeatother
\usepackage{titlesec}
\usepackage{fancyhdr}
\usepackage[unicode,hidelinks,bookmarksnumbered]{hyperref}

\linespread{1.12}
\setlength{\parindent}{0pt}
\setlength{\parskip}{0.55em}
\exhyphenpenalty=10000          % never break a numeric range after its dash
\emergencystretch=4em
\setcounter{tocdepth}{2}
\setcounter{secnumdepth}{2}
\renewcommand{\topfraction}{0.9}
\renewcommand{\bottomfraction}{0.5}
\renewcommand{\textfraction}{0.07}
\renewcommand{\floatpagefraction}{0.85}
\setcounter{topnumber}{3}
\setlist{itemsep=0.25em,topsep=0.3em}

% ---- colours (same tokens as the figures) -------------------------------------------------
\definecolor{ink}{HTML}{0B0B0B}
\definecolor{ink2}{HTML}{52514E}
\definecolor{rule}{HTML}{E4E3DF}
\definecolor{accent}{HTML}{2A78D6}
\definecolor{aqua}{HTML}{1BAF7A}
\definecolor{goodc}{HTML}{1F7A3A}
\definecolor{badc}{HTML}{B42318}
\definecolor{warnc}{HTML}{A15C00}
\definecolor{neuc}{HTML}{52514E}

% ---- headings, header/footer --------------------------------------------------------------
\titleformat{\section}{\LARGE\bfseries}{\thesection.}{0.5em}{}[\vspace{2pt}{\color{ink}\titlerule[1.2pt]}]
\titlespacing*{\section}{0pt}{0pt}{1em}
\titleformat{\subsection}{\Large\bfseries}{\thesubsection}{0.6em}{}
\titlespacing*{\subsection}{0pt}{1.4em}{0.5em}
\titleformat{\subsubsection}[runin]{\bfseries}{}{0pt}{}[.]
\titlespacing*{\subsubsection}{0pt}{0.9em}{0.5em}
\newcommand{\chapterbreak}{\clearpage}
\let\oldsection\section
\renewcommand{\section}{\chapterbreak\oldsection}

\pagestyle{fancy}
\fancyhf{}
\fancyhead[R]{\small\color{ink2}טופולוגיה של עולם סמוי בייצוגי \en{Transformer}}
\fancyhead[L]{\small\color{ink2}\nouppercase{\leftmark}}
\fancyfoot[C]{\small\thepage}
\renewcommand{\headrulewidth}{0.4pt}
\renewcommand{\sectionmark}[1]{\markboth{#1}{}}

% ---- boxes and verdict chips --------------------------------------------------------------
\tcbset{hebox/.style={enhanced,breakable,arc=0pt,outer arc=0pt,boxrule=0pt,rightrule=3pt,
  leftrule=0pt,left=10pt,right=10pt,top=6pt,bottom=6pt,before skip=0.9em,after skip=0.9em,
  fonttitle=\bfseries,coltitle=ink,attach title to upper={\par\smallskip}}}
\newtcolorbox{keybox}[1][]{hebox,colback=accent!6,colframe=accent,title={#1}}
\newtcolorbox{meaningbox}[1][משמעות לאור שאלות המחקר]{hebox,colback=aqua!7,colframe=aqua,title={#1}}
\newtcolorbox{caveat}[1][הסתייגות]{hebox,colback=warnc!6,colframe=warnc,title={#1}}
\newcommand{\chip}[2]{\tcbox[on line,arc=6pt,boxrule=0.5pt,colframe=#1,colback=#1!8,
  left=3pt,right=3pt,top=0pt,bottom=0pt,boxsep=1pt,fontupper=\footnotesize\bfseries\color{#1}]{#2}}
\newcommand{\good}[1]{\chip{goodc}{#1}}
\newcommand{\bad}[1]{\chip{badc}{#1}}
\newcommand{\warn}[1]{\chip{warnc}{#1}}
\newcommand{\neu}[1]{\chip{neuc}{#1}}

\newcolumntype{C}{>{\centering\arraybackslash}X}
\newcolumntype{Y}{>{\raggedleft\arraybackslash}X}
\renewcommand{\arraystretch}{1.25}
\newcommand{\fig}[3][\linewidth]{%
  \begin{figure}[!ht]\centering\includegraphics[width=#1]{figs/#2}\caption{#3}\label{fig:#2}\end{figure}}

% ---- notation -----------------------------------------------------------------------------
\newcommand{\Sc}[3][]{S^{\mathrm{#2}}_{\mathrm{#3}#1}}
\newcommand{\bel}[1]{b^{(#1)}}
\newcommand{\GC}{G_{C}}
\newcommand{\KL}{\mathrm{KL}}
\newcommand{\dgeom}{\Delta_{\mathrm{geom}}}

\hypersetup{pdftitle={טופולוגיה של עולם סמוי בייצוגים של Transformer},pdfauthor={topo-pilot}}

\begin{document}

% =========================================================================================
\begin{titlepage}
\thispagestyle{empty}
\vspace*{5cm}
{\large\color{ink2} דוח מחקר מסכם · פרויקט \en{topo-pilot}\par}
\vspace{0.6cm}
{\fontsize{34}{40}\selectfont\bfseries טופולוגיה של עולם סמוי\par בייצוגים של \en{Transformer}\par}
\vspace{0.6cm}
{\Large\color{ink2} שאלות המחקר, זרימת העבודה, המתודולוגיה של כל שלב,\par תוצאות הניסויים ומשמעותן\par}
\vspace{1.2cm}
{\color{rule}\rule{\linewidth}{0.6pt}}\par
{\color{ink2} גרסה 2.0 · 15 באוקטובר 2026 · סודר ב־\en{LuaLaTeX}\par
מבוסס על המאגר \en{topo-pilot} (ענף \en{main}), על הדוחות \code{reports/E0-E4}, על קובצי הרישום המוקדם ועל רשומות מחקרי ההמשך.\par}
\vspace{0.8cm}
\good{E4 הושג}\quad \bad{שער השחזור G-H נכשל ומתועד}\quad \neu{אין טענה סיבתית}
\end{titlepage}

\pagenumbering{roman}
\setcounter{page}{2}
\tableofcontents
\vspace{1em}
\begin{keybox}[איך לקרוא את הדוח]
הדוח בנוי בשלוש שכבות. סעיפים 1–3 נותנים תמונה כוללת: מה שאלנו, איך זורם המחקר ואיך מתקבלות ההכרעות. סעיף 4 מפרט את המתודולוגיה של כל שלב, כולל הנימוק לכל בחירה. סעיפים 5–7 מציגים את התוצאות, ולצד כל תוצאה תיבה ירוקה שמסבירה את משמעותה לאור שאלות המחקר. תיבות כתומות מסמנות הסתייגויות. ההכרעות המחייבות מנוסחות באנגלית ב־\code{README.md}; הדוח הזה הוא הסבר בעברית, ואינו משנה אף הכרעה.
\end{keybox}

\clearpage
\pagenumbering{arabic}

% =========================================================================================
\section{תקציר מנהלים}\label{sec:summary}

המחקר בודק שאלה בסיסית על מודלי שפה: כאשר מאמנים \en{Transformer} לחזות את הסמל הבא ברצף, האם הייצוג הפנימי שלו לומד משהו על \textbf{צורת העולם} שיצר את הרצף? כדי שאפשר יהיה לענות בלי ניחושים, בנינו עולם סינתטי שבו הכול ידוע: שני עולמות בגודל $12\times12$, \textbf{גליל} (שיש בו ״חור״ אחד, $\beta_1=1$) ו\textbf{ריבוע} (שאין בו חור, $\beta_1=0$). סוכן נע בהם בהליכה מקרית, ובכל צעד פולט סמל רועש מתוך 36 סמלים. המודל רואה רק את הסמלים. אחרי האימון בדקנו, בעזרת \textbf{הומולוגיה מתמידה} (\en{Persistent Homology, PH}), אם ענן האקטיבציות של המודל מכיל את המעגל של הגליל, והשווינו אותו למנבא בייסיאני מדויק שמחושב מהמחולל.

העבודה התקדמה בשלבים \en{S0--S8}, עם חמש נקודות יציאה (\en{E0--E4}). כל כלל הכרעה הוקפא ברישום מוקדם לפני שנראתה התוצאה הרלוונטית. אומנו 17 מודלים על מעבד בלבד. אחרי \en{E4} בוצעו תשעה מחקרי המשך, ורובם נרשמו מראש.

\begin{keybox}[חמשת הממצאים העיקריים]
\begin{enumerate}[leftmargin=*]
\item \textbf{תחת פליטה חלקה (O), האימון יוצר אות $H_1$ שמבחין בין גליל לריבוע} (F1 נתמך): $D>0$ ב־8 מתוך 8 לוחות בכל אחת משלוש החזרות, והחציון (0.158–0.176) גבוה ממעטפת הרעש (0.084–0.097). מחלקת ה־$H_1$ הבולטת היא אכן המעגל המחזורי של העולם: התאמה של 0.88–0.89 לזווית הסמויה, לעומת 0.92 באמונה המדויקת.
\item \textbf{הגאומטריה של הייצוג עוקבת אחרי האמונה שנבנתה מכל ההיסטוריה}, ולא רק אחרי הסמל האחרון (H2: $\dgeom>0$ ב־48 מתוך 48 לוחות, גם בתוך זוגות היסטוריות עם אותו סמל אחרון).
\item \textbf{תחת פליטה מעורבבת (P) אין ניגוד $H_1$}, אבל זו תוצאה לא־אינפורמטיבית מראש, כי גם האמונה המדויקת אינה מבחינה שם בין העולמות. ובכל זאת, הייצוג מתיישר עם גאומטריית האמונה הבייסיאנית גם תחת P (H3$'$ נתמך).
\item \textbf{השערת התזמון (H4) נדחתה:} ציון ה־$H_1$ עולה כבר בצעד 50, לפני שהמודל מתחיל להשתמש בהיסטוריה (צעד 100), ובשני העולמות. מקור המבנה, שילוב היסטוריה או סטטיסטיקה ברמת הסמל, נשאר \textbf{פתוח}.
\item \textbf{אין ראיה סיבתית.} כל הממצאים ייצוגיים. \en{probes} לינאריים מפענחים את האמונה מהאקטיבציות באותה רמה, ולא יותר, מאשר מהתפלגות הסמל הבא המדויקת. ניסיון ההתערבות (\en{X4-F}) נעצר בשער היתכנות.
\end{enumerate}
\end{keybox}

\begin{table}[H]\small
\begin{tabularx}{\linewidth}{@{}l>{\raggedleft\arraybackslash}p{0.33\linewidth}l Y@{}}
\toprule
\textbf{השערה} & \textbf{מה היא אומרת} & \textbf{הכרעה} & \textbf{משמעות בקצרה} \\
\midrule
\textbf{F1} & האימון מגביר אות $H_1$ שמבחין בין גליל לריבוע (תחת O) & \good{נתמך} & הייצוג נושא את הטופולוגיה של העולם \\
\textbf{H2} & הגאומטריה עוקבת אחרי אמונת 32 הסמלים יותר מאחרי הסמל האחרון & \good{עקבי בסימן} & יש שימוש בהיסטוריה, לפחות בגאומטריה \\
\textbf{H3} & האות שורד גם תחת P & \bad{נדחה} & לא אינפורמטיבי; לא ניתן להכריע בכיוון הזה \\
\textbf{H3$'$} & תחת P הגאומטריה מתיישרת עם האמונה, מעבר לאתחול ולבסיסי הערבוב & \good{נתמך} & התיישרות עם אמונה גם בלי ניגוד $H_1$ \\
\textbf{H4} & $S$ עולה רק אחרי שהמודל מתחיל להשתמש בהיסטוריה & \bad{נדחה} & העלייה המוקדמת אינה ייחודית לגליל \\
\textbf{H0} & בפליטות זהות לא נלמד דבר & \good{מתקיים} & אין זליגת מידע; הצינור תקין \\
\textbf{G-H} & שחזור נקי מאפס & \bad{נכשל} & 307/307 תוצרים משותפים זהים ביט־לביט; רשומה היסטורית אחת חסרה \\
\bottomrule
\end{tabularx}
\caption{ההשערות והכרעותיהן במבט אחד.}\label{tab:verdicts}
\end{table}

% =========================================================================================
\section{שאלות המחקר וההשערות}\label{sec:questions}

\subsection{השאלה המרכזית}
כשמאמנים מודל לחזות את הסמל הבא על רצפים שנוצרו מעולם בעל טופולוגיה ידועה, עולות שלוש שאלות, וכל אחת בנויה על קודמתה:
\begin{enumerate}[leftmargin=*]
\item \textbf{שאלת הקיום (ש1):} האם ב־\en{residual stream} של המודל נוצר מבנה שמשקף את הטופולוגיה של העולם הסמוי?
\item \textbf{שאלת המקור (ש2):} אם כן, מאין הוא בא? יש שלושה מקורות אפשריים: (א) \textbf{גאומטריית התצפיות}, כלומר הדמיון בין סמלים שנפלטים ממקומות קרובים; (ב) \textbf{האמונה המצטברת} על מיקום הסוכן, שנבנית מכל ההיסטוריה; (ג) רק \textbf{מה שנחוץ לחיזוי הסמל הבא}, בלי עודף.
\item \textbf{שאלת השימוש (ש3):} האם המודל \textbf{משתמש} במבנה הזה כדי לחזות, או שהוא רק תוצר לוואי?
\end{enumerate}
התוכנית המקורית ענתה במפורש רק על ש1 ועל חלק מ־ש2. ש3 דורשת התערבות באקטיבציות, והוגדרה מראש כהרחבה. ההבחנה הזו חשובה לקריאת כל התוצאות: \textbf{הפרויקט מזהה מבנה ומודד את המתאם שלו עם החיזוי; הוא אינו מראה שהמודל צריך את המבנה.}

\subsection{ההשערות הניתנות להפרכה}
כל שאלה תורגמה להשערה עם מדד מוגדר ועם תנאי הפרכה שנקבע מראש (טבלה~\ref{tab:hyp}). ההשערה הראשית היא אחת בלבד (F1), כדי לא לפזר את כוח ההכרעה על פני מבחנים רבים. שאר ההשערות נועדו לפרש אותה.

\begin{table}[H]\small
\begin{tabularx}{\linewidth}{@{}l Y >{\raggedleft\arraybackslash}p{0.2\linewidth} Y l@{}}
\toprule
\textbf{מזהה} & \textbf{השערה} & \textbf{מדד} & \textbf{מופרכת או מוגבלת כאשר} & \textbf{שאלה}\\
\midrule
\textbf{F1} & האימון מגביר אות $H_1$ מתמיד שמבחין בין גליל לריבוע & $D_{r,i}$ & תנאי כלשהו בכלל ההכרעה נכשל, כשכל השערים עברו & ש1\\
\textbf{H2} & הגאומטריה עוקבת אחרי אמונת 32 הסמלים יותר מאחרי אמונת הסמל האחרון & $\dgeom$, גם בזוגות עם אותו סמל אחרון & $\dgeom\le 0$, או שהאפקט נעלם בהתניה על הסמל האחרון & ש2(ב)\\
\textbf{H3} & תחת P עדיין $D>0$, כלומר המבנה נלמד מהדינמיקה ולא רק מגאומטריית הפליטה & כלל F1 תחת P & הכלל נכשל תחת P אף שהמודל עבר את שער הלמידה & ש2(א)\\
\textbf{H4} & $S$ עולה רק אחרי שהמודל מתחיל להשתמש בהיסטוריה & עקומות על 12 נקודות בדיקה & $S$ עולה לפני השימוש בהיסטוריה & ש2(ב)\\
\textbf{H0} & בפליטות זהות לא ניתן ללמוד דבר על המצב & \en{NLL} מול $H(W)$ & \en{NLL} נמוך מ־$H(W)$: זליגה או באג & תקינות\\
\bottomrule
\end{tabularx}
\caption{ההשערות שנרשמו מראש, המדדים ותנאי ההפרכה.}\label{tab:hyp}
\end{table}

\subsection{זיהוי, שימוש והכרח}
שלושת המושגים האלה נשמרו נפרדים לאורך כל העבודה. \textbf{זיהוי} פירושו שיש בענן האקטיבציות מבנה שאפשר למדוד, למשל פס $H_1$ ארוך. \textbf{שימוש} פירושו שהמודל מסתמך על המבנה בחישוב התחזית, ואפשר להראות זאת רק בהתערבות: משנים את המבנה ובודקים אם החיזוי נפגע. \textbf{הכרח} פירושו שאי אפשר לחזות היטב בלי המבנה, וזו טענה על המשימה ולא על המודל. מרחבי העתיד $F_1,F_2$ מראים אילו הבחנות נחוצות לחיזוי עד אופק 2, ולא יותר.

בנוסף, הפרויקט אינו מניח שאקטיבציות או אמונות הן יריעות. הומולוגיה של קומפלקס \en{Vietoris--Rips} היא תכונה של ענן נקודות ושל מטריקה בסקאלה נתונה. היא אינה מוכיחה הומאומורפיזם, ו־$S$ גבוה אינו שקול ל־$\beta_1=1$.

\subsection{עקרונות מתודולוגיים מנחים}
\begin{itemize}[leftmargin=*]
\item \textbf{רישום מוקדם.} ההשערות, המדדים, הספים וכללי ההכרעה נכתבו ונשמרו ב־\en{git} לפני שנראו התוצאות הרלוונטיות (\code{PREREGISTRATION.md} בסוף S2; \code{PREREGISTRATION_X4.md} לניתוחים המשניים). שינוי בהחלטה סגורה מחייב עצירה ואישור הבעלים (נקודה H3).
\item \textbf{נקודות יציאה.} כל שלב מסתיים בתוצר תקין ובר־הצגה. אם העבודה נעצרת, מה שכבר נעשה עומד בפני עצמו, ומה שלא נעשה מסומן במפורש ״לא הורץ״.
\item \textbf{שערי תוקף.} לפני שמפרשים תוצאה, בודקים שהמדידה עצמה עובדת (\en{G-A} עד \en{G-H}). אם המודל לא למד, היעדר מבנה פירושו ״המודל לא למד״, ולא ממצא שלילי.
\item \textbf{עקביות בסימן במקום \en{p-value}.} נקודות בענן אחד אינן תצפיות בלתי תלויות, ולכן אין מבחני מובהקות על זוגות. ההכרעה מבוססת על כיוון עקבי בשלוש חזרות ובשמונה לוחות זרים, ועל מעטפת רעש שנגזרת מהנתונים.
\item \textbf{בקרות מותאמות.} כל טענה מלווה בבקרה שמנטרלת הסבר חלופי: אתחול (R0), סמל אחרון (R1), ערבוב פשוט (R2, R3), רצפה גאוסית, פליטה מעורבבת (P), פליטה זהה (N) ומצב גלוי (U).
\item \textbf{הפרדה בין מאושר לחקרני.} כל תצפית שנעשתה אחרי חשיפת התוצאות נכתבת בקובץ נפרד (\code{exploratory.md}, \code{interpretation.md}), ואינה משנה הכרעה.
\item \textbf{בדיקות בלתי תלויות.} במודולים הקריטיים (עולם, טופולוגיה, אורקל, כלל F1) את הבדיקות כתב סוכן אחר מזה שכתב את הקוד, לפי תוכנית בדיקות כתובה.
\end{itemize}

% =========================================================================================
\section{זרימת המחקר}\label{sec:flow}

\subsection{זרימת הנתונים והניתוח}
איור~\ref{fig:pipeline} מראה את מסלול המידע מהמחולל ועד ההכרעה. העיקרון המרכזי הוא \textbf{שני ענפים מקבילים על אותן היסטוריות}. בענף הלמידה, \en{Transformer} מאומן מאפס על הסמלים בלבד, ומהאקטיבציות שלו נבנה ענן נקודות. בענף האורקל, מנבא בייסיאני מדויק מחשב מאותן היסטוריות את האמונה על המצב הסמוי ואת התפלגויות העתיד. שני הענפים נפגשים בלוחות דגימה משותפים, ושם עוברים בדיוק את אותו צינור ניתוח. כך אפשר לשאול אם הגאומטריה שהמודל בנה דומה לגאומטריה של מה שמנבא אידאלי היה יודע.

\fig[0.92\linewidth]{pipeline}{זרימת הנתונים והניתוח. המחולל (שורה עליונה) ידוע במלואו. הקואורדינטות של המצבים זמינות רק למחולל ולניתוח, ולעולם לא למודל. שני הענפים נפגשים ב־8 לוחות זרים של 512 היסטוריות, והניתוחים כולם רצים על אותן היסטוריות בכל מרחב ובכל נקודת בדיקה.}

\subsection{שלבים, שערים ונקודות יציאה}
המחקר חולק לתשעה שלבי הרצה (איור~\ref{fig:stages}). לכל שלב יש שער אוטומטי שחייב לעבור לפני שממשיכים, ובארבע נקודות העבודה נעצרת לאישור הבעלים: H1 אחרי הכיול והרישום המוקדם; H2 לפני האימונים, עם הערכת זמן; H3 לפני כל שינוי בהחלטה סגורה; H4 לפני כתיבת הממצא. חמש נקודות יציאה מגדירות מה מותר לטעון בכל עצירה (טבלה~\ref{tab:exits}). הבנייה עצמה נעשתה ב״גלים״ מקבילים (\en{W0--W5}), שכל אחד מהם מאפשר את שלב ההרצה המתאים.

\fig[0.86\linewidth]{stages}{שלבי המחקר. שלבים \en{S0--S8} מימין, השער האוטומטי של כל שלב, נקודות האישור של הבעלים (מעוינים) ונקודות היציאה. השער \en{G-H} נכשל בשלב S8 (סעיף~\ref{sec:repro}).}

\begin{table}[H]\small
\begin{tabularx}{\linewidth}{@{}l l Y c@{}}
\toprule
\textbf{יציאה} & \textbf{אחרי} & \textbf{מה מותר לטעון} & \textbf{ריצות אימון}\\
\midrule
\textbf{E0} & S1 & צינור ה־PH מזהה את הטופולוגיה הידועה; מחזורי גרף אינם חורי משטח & 0\\
\textbf{E1} & S2 & ההבדלים הטופולוגיים קיימים בעיקרון, בלי מודל & 0\\
\textbf{E2} & S4 & הכרעה ב־F1 (נתמך, נדחה או לא מכריע) ותוצאת H2 & 7\\
\textbf{E3} & S5 + S6 & H3 ו־H4 & 13\\
\textbf{E4} & S7 + S8 & פרשנות מחוזקת: בקרות, \en{probes}, קואורדינטות מעגליות, רגישות, שחזור & 17\\
\bottomrule
\end{tabularx}
\caption{נקודות היציאה ומה מותר לטעון בכל אחת.}\label{tab:exits}
\end{table}

\subsection{לוגיקת ההכרעה}
ההכרעה בהשערה הראשית נעשית בכלל דטרמיניסטי, שנכתב לפני שנאסף נתון כלשהו מהמודלים (איור~\ref{fig:decision}). יש לו שלוש תוצאות אפשריות, ולא שתיים: ״לא מכריע״ שמור למקרה שבו שער תוקף נכשל, כדי שכשל טכני לא ייקרא בטעות כממצא שלילי.

\fig[0.9\linewidth]{decision}{כלל ההכרעה F1. התנאי הראשון דורש עקביות בכיוון (לפחות 7 מתוך 8 לוחות בכל חזרה). התנאי השני דורש שהאפקט גדול מהשונות בין לוחות בתוך אותה ריצה. אותו כלל, בלי שינוי, הוחל גם תחת P, בכל אחת מנקודות המבט \en{V0--V5} ובכל וריאנט רגישות.}

לצד הכלל נכתבה מראש \textbf{טבלת פירוש} (טבלה~\ref{tab:interp}), שקובעת מה מותר להסיק מכל צירוף תוצאות. כך נמנע מצב שבו הפרשנות מותאמת לתוצאה אחרי מעשה.

\begin{table}[H]\small
\begin{tabularx}{\linewidth}{@{}>{\raggedleft\arraybackslash}p{0.36\linewidth} Y@{}}
\toprule
\textbf{צירוף תוצאות} & \textbf{מסקנה מותרת}\\
\midrule
F1, H2 ו־H3 נתמכים & אות $H_1$ שנלמד מההיסטוריה ומהדינמיקה ומבחין בין העולמות; עדיין לא הוכחה לשימוש\\
F1 נתמך תחת O ולא תחת P & המבנה ככל הנראה עובר משדה הפליטה, \textbf{אלא אם} גם האמונה המדויקת אינה מבחינה תחת P (וזה מה שקרה, סעיף~\ref{sec:e1})\\
F1 נתמך, H2 נדחה & המבנה מוסבר על ידי הסמל האחרון\\
F1 נתמך, ו־R2 או R3 מראים $D$ באותו סדר גודל & האות מוסבר בערבוב סמלים פשוט; המסקנה נשענת אז על P\\
F1 נדחה, המודל למד, והאמונות מבחינות & חיזוי טוב בלי אות $H_1$ גולמי במדידה הזו; ייתכן קידוד בגאומטריה אחרת\\
המרחבים הלא־נלמדים אינם מבחינים & אין חתימה צפויה; השאלה אינה ניתנת להכרעה במערך הזה\\
הצינור אינו מזהה את הגליל הידוע & כשל מדידה; לא מפרשים אקטיבציות\\
\bottomrule
\end{tabularx}
\caption{טבלת הפירוש שנקבעה מראש.}\label{tab:interp}
\end{table}

% =========================================================================================
\section{המתודולוגיה המחקרית לפי שלב}\label{sec:method}
סעיף זה מסביר כל שלב: מה נעשה בו, למה דווקא כך, ואיזה שער מבטיח שאפשר לסמוך על התוצר שלו. הבחירות מוסברות כאן כי רובן נועדו לסגור הסבר חלופי מסוים, והקורא צריך לדעת איזה.

\subsection{שלב S0: תשתית, חוזים ובקרת איכות}
לפני כל ניסוי נבנתה שכבת \textbf{חוזים}: טיפוסי נתונים קבועים לכל תוצר (מפרט עולם, מפרט פליטה, מערך תצפיות, פלטי אורקל, נקודת בדיקה, מערך אקטיבציות, רשומות מדדים, דוחות שערים ועוד). כל תוצר נשמר ב\textbf{מאגר תוצרים} עם מזהה שנגזר מתוכנו ומהוריו, ועם \en{checksum} של המטען. מכאן נובעות שתי תכונות. ראשית, כל מספר בדוח ניתן לעקיבה לתוצר שיצר אותו. שנית, שחזור נבדק בהשוואת מזהים ו־\en{checksums}, ולא בהשוואת מספרים מעוגלים.

הקוד חולק למודולים עם ממשקים ציבוריים בלבד, ונבנה בגלים מקבילים. במודולים הקריטיים (\code{world}, \code{topology}, \code{oracle}, \code{analysis}) כותב הבדיקות היה סוכן אחר מזה שמימש, ועבד מתוכנית בדיקות כתובה, כדי שבדיקה לא ״תתאים את עצמה״ לבאג. השער של S0: חבילת הבדיקות של החוזים עוברת, ובדיקת התלויות בין המודולים נקייה. הסביבה ננעלה (\code{requirements.lock}) והכול רץ על מעבד, בלי \en{GPU}.

\subsection{שלב S1: העולמות ואמת המידה הטופולוגית}
\subsubsection{העולמות}
שני העולמות הם סריג של $12\times12$ מצבים עם שכנות של 4 (איור~\ref{fig:world}). בריבוע אין הדבקה. בגליל ציר $x$ מחזורי: יציאה מהעמודה הימנית מובילה לעמודה השמאלית. הבחירה בגודל זהה ($N=144$) בשני העולמות מבטלת הבדל בכמות המצבים. עולם ״רצועה״ של $16\times4$, שהופיע בטיוטה מוקדמת, הוחלף, כי חצי ממצביו ישבו על השפה.

\fig[0.9\linewidth]{world}{שני העולמות. 144 מצבים בכל עולם. הקשתות המקווקוות בגליל מסמנות את ההדבקה בין שתי העמודות הקיצוניות. המעוינים הם 36 מרכזי הפליטה (סעיף~\ref{sec:s2}); הם אינם מצבים.}

\subsubsection{אמת המידה: גרף אינו משטח}
הטופולוגיה האמיתית חושבה מ\textbf{קומפלקס קובייתי} מפורש: קודקודים, צלעות ופאות ריבועיות מלאות. מספרי בטי חושבו מדרגות מטריצות השפה $\partial_1,\partial_2$ מעל השדה $\mathbb{F}_2$, עם בדיקה ש־$\partial_1\partial_2=0$. ההבחנה חשובה: לגרף המעברים של הריבוע יש 121 מחזורים בלתי תלויים, אחד סביב כל פאה, ולגליל 133. אבל בקומפלקס הפאות מלאות, ולכן הן אינן חורים. ״מחזור בגרף״ ו״חור במשטח״ הם דברים שונים, וכל הניתוח מכוון לשני.

\subsubsection{הדינמיקה}
הסוכן נע בהליכה מקרית \textbf{עצלה}: בהסתברות 0.2 הוא נשאר במקום, ובהסתברות 0.2 הוא זז לכל אחד מארבעת הכיוונים. ניסיון לצאת מהשפה הופך להישארות. בכתיבה כללית, לכל זוג שכנים $i\sim j$:
\begin{equation}
T_{ij}=\frac{1-p_{\mathrm{stay}}}{D_{\max}},\qquad
T_{ii}=1-\deg(i)\,\frac{1-p_{\mathrm{stay}}}{D_{\max}},\qquad
p_{\mathrm{stay}}=0.2,\; D_{\max}=4 .
\end{equation}
המטריצה סימטרית ודו־סטוכסטית, ולכן ההתפלגות הסטציונרית $\pi$ אחידה. לולאות עצמיות מבטיחות שהשרשרת אינה מחזורית. ההבדלים שנשארים בין העולמות (שיעור ההישארות בשפה, הקוטר, הפער הספקטרלי) נמדדו ודווחו: זו השוואה בין שני תהליכים, ולא התערבות בטופולוגיה בלבד.

\subsubsection{המטריקה ושער \en{G-A}}
המטריקה הראשית היא המרחק האינטרינזי השטוח $d_M$, עם מינימום מחזורי בציר $x$ בגליל:
\begin{equation}
d_M\big((x,y),(x',y')\big)=\sqrt{\min\!\big(|x-x'|,\,12-|x-x'|\big)^2+(y-y')^2},
\end{equation}
בגליל, ובלי המינימום בריבוע. מטריקת הרגישות היא מרחק בגרף. השער \en{G-A} בודק שהצינור כולו עובד על אמת ידועה, לפני שמפעילים אותו על אקטיבציות. הוא דורש שלושה דברים: (1) מספרי בטי נכונים בשני העולמות; (2) \en{VR} על $d_M$ נותן פס $H_1$ ארוך לגליל בלבד, עם ספים שנקבעו מראש ($S\ge0.2$ לגליל, $S\le0.1$ לריבוע); (3) ענני מעגל ודיסק סינתטיים, משוכנים ב־$\mathbb{R}^{128}$ עם רעש ו־512 נקודות, מסווגים נכון ב־19 מתוך 20 הגרלות לפחות.

\subsection{שלב S2: פליטות, מנבא מדויק וכיול}\label{sec:s2}
\subsubsection{ארבעה תנאי פליטה}
כל מצב פולט סמל לפי מטריצת פליטה $E$. נבנו ארבעה תנאים, וכל אחד עונה על שאלה אחרת (טבלה~\ref{tab:emissions}).

\begin{table}[H]\small
\begin{tabularx}{\linewidth}{@{}l >{\raggedleft\arraybackslash}p{0.42\linewidth} Y l@{}}
\toprule
\textbf{קוד} & \textbf{הגדרה} & \textbf{תפקיד} & \textbf{שלב}\\
\midrule
\textbf{O} & 36 מרכזים ברשת $6\times6$, $c_{ab}=(0.5+2a,\,0.5+2b)$; $E_{zj}\propto\exp\!\big(-d_M(z,c_j)^2/2\sigma^2\big)$; שמות הסמלים מעורבבים & הניסוי המרכזי: אי־ודאות רציפה שההיסטוריה מצמצמת & S4\\
\textbf{P} & פרמוטציה לפי צירים של O: $E_P[(x,y)]=E_O[(\sigma_x(x),\sigma_y(y))]$ & בקרה מותאמת: אותן שורות, אותה אנטרופיה, אותה מכפלתיות; החלקות נהרסת & S5\\
\textbf{N} & $E_{zw}=1/36$ & בקרה שלילית: אין מה ללמוד על המצב & S4, S7\\
\textbf{U} & סמל ייחודי לכל מצב, $V=144$ & בקרת תקינות: המצב גלוי, לומדים רק דינמיקה & S7\\
\bottomrule
\end{tabularx}
\caption{תנאי הפליטה.}\label{tab:emissions}
\end{table}

\textbf{למה P בנויה לפי צירים.} בסקירת הספרות התברר שהפליטה O מתפרקת במדויק למכפלה:
\begin{equation}
E_O\big[(x,y),(a,b)\big]=E^{x}[x,a]\;E^{y}[y,b].
\end{equation}
כל סמל הוא בעצם זוג של קריאות רועשות, אחת לכל ציר. מחקרים קודמים מראים שמודלים נוטים לייצג גורמים כמעט בלתי תלויים בתת־מרחבים נפרדים. פרמוטציה כללית של 144 המצבים הייתה שוברת גם את החלקות וגם את המכפלתיות, ומערבבת את שתי ההשפעות. פרמוטציה לפי צירים שוברת רק את החלקות, ולכן ההבדל בין O ל־P מבודד את תרומת גאומטריית התצפיות (איור~\ref{fig:emission}).

\fig[0.9\linewidth]{emission}{שדה הפליטה של סמל אחד. תחת O (מימין) ההסתברות לפלוט את הסמל גבוהה באזור רציף אחד, ובגליל היא נמשכת דרך התפר לצד השני. תחת P (משמאל) אותם ערכים בדיוק מפוזרים בין שורות ועמודות שאינן שכנות. העמודות והשורות נשמרות כיחידות, ולכן נשמרות גם המכפלתיות וגם התפלגות השורות.}

\subsubsection{המנבא הבייסיאני המדויק}
הרצפים נוצרים כך: מצב התחלתי מ־$\pi$, ובכל צעד פליטה ואחריה מעבר, עד 33 סמלים. המודל מקבל $[\mathrm{BOS},w_1,\dots,w_{32}]$ וחוזה בכל עמדה את הסמל הבא. כי כל מסלול מתחיל מ־$\pi$, אפשר לחשב בדיוק את מה שמנבא אידאלי היה יודע באותה עמדה, בסינון בייסיאני (וקטורי שורה, \en{float64}):
\begin{equation}
a_1=\pi,\qquad
b_t=\frac{a_t\odot E_{:,w_t}}{a_t\,E_{:,w_t}},\qquad
a_{t+1}=b_t\,T,\qquad
q_t=b_t\,T\,E .
\end{equation}
מזה נגזרים מרחבי הייחוס: $\bel{32}$, האמונה על המצב אחרי כל 32 הסמלים; $\bel{1}=r_{w_{32}}$, האמונה מהסמל האחרון בלבד, כאשר $r_w=\pi\odot E_{:,w}/(\pi E_{:,w})$; $F_1=b\,TE$, התפלגות הסמל הבא; ו־$F_2$, ההתפלגות המשותפת של שני הסמלים הבאים:
\begin{equation}
F_2[i,j]=b\,T\,\mathrm{diag}(E_{:,i})\,T\,E_{:,j}.
\end{equation}
ההשוואה בין $\bel{32}$ ל־$\bel{1}$ מאפשרת לשאול מה ההיסטוריה מוסיפה, בלי לחשוף את המצב האמיתי. ראיה אפסית מעלה שגיאה מפורשת, בלי תיקון שקט באפסילון.

מדד \textbf{תועלת ההקשר} מודד כמה ההיסטוריה המלאה משפרת את החיזוי מעבר לסמל האחרון, בעמדה 32:
\begin{equation}
\GC=\overline{\KL\big(q_C\,\|\,q_{\mathrm{last}}\big)},\qquad q_{\mathrm{last}}=r_{w_{32}}\,T\,E .
\end{equation}
אם הוא קטן, אין סיבה לצפות שהמודל יבנה אמונה מההיסטוריה.

\subsubsection{כיול $\sigma$}
רוחב הפליטה $\sigma$ נבחר \textbf{מהמחולל בלבד}, לפני כל אימון, לפי סדר קבוע ($1.25$, אחר כך $0.75$, אחר כך $2.0$). נבחר הערך הראשון שעומד בשלושה תנאים בשני העולמות: האנטרופיה הממוצעת של $\bel{1}$ לפחות $\log 2$ (המצב אינו כמעט גלוי); $\GC\ge0.10$ \en{nats} (להיסטוריה יש ערך ממשי); וכל סמל מופיע בהסתברות $p(w)\ge0.5/V$. הערך הראשון, $\sigma=1.25$, עמד בכל התנאים ($\GC=0.265$ בגליל ו־$0.236$ בריבוע), ולכן לא נדרש מסלול הנפילה.

\subsubsection{השערים \en{G-B}, \en{G-C}, \en{G-D} והרישום המוקדם}
\en{G-B} בודק שהאורקל תואם ספירה מלאה בכוח גס על מקרים קטנים, ושבקרות U ו־N מתנהגות כמצופה. \en{G-C} בודק שה־$\sigma$ נבחר לפי הכלל. \en{G-D} בודק שאין בנתונים הציבוריים שדה שחושף את המצב הסמוי, ושבדיקות הזליגה עוברות. בסוף השלב חושבה טבלת ההבחנה של המרחבים הלא־נלמדים (\en{E1}, סעיף~\ref{sec:e1}), ונכתב הרישום המוקדם, עם \en{hash} של הקונפיגורציה.

\subsection{שלב S3: המודל, האימון ושער הלמידה}
המודל הוא \en{Transformer} מסוג \en{decoder-only} עם \en{pre-LN}: 4 בלוקים, ממד 128, 4 ראשים, \en{MLP} ברוחב 512, \en{embedding} מיקום נלמד, \en{LayerNorm} סופי ו־\en{unembedding} לא קשור; כ־0.8 מיליון פרמטרים. האימון: \en{AdamW} עם קצב למידה $10^{-3}$, 200 צעדי חימום ודעיכת \en{cosine} עד 5000 צעדים, אצווה של 128, כ־21 מיליון מטרות לריצה. לכל חזרה 131,072 מסלולי אימון, 4,096 לפיתוח ו־4,096 למבחן. שלושה זרעים (1729, 2718, 3141) משנים יחד את הנתונים, האתחול, שמות הסמלים וסדר האצוות.

\textbf{נקודות בדיקה בפריסה לוגריתמית} נשמרות תמיד: 0, 50, 100, 200, 400, 700, 1000, 1500, 2000, 3000, 4000, 5000. הניתוח הראשי נעשה בצעד 5000 ובצעד 0, בלי לבחור את נקודת הבדיקה ״היפה״ ביותר.

\textbf{שער הלמידה.} המדד הראשי לחיזוי הוא המרחק מהמנבא המדויק, שמחושב במדויק לכל היסטוריה, וממנו נגזר החלק מתועלת ההקשר שהמודל תפס:
\begin{equation}
\mathrm{regret}=\overline{\KL\big(q_C\,\|\,p_{\mathrm{model}}\big)},\qquad
\mathrm{captured}=1-\frac{\mathrm{regret}}{\GC}\;\ge\;0.8 .
\end{equation}
השער \en{G-E} הורץ על ריצת ניסיון אחת (גליל, O, זרע 1729), ובדק גם שהחישוב ייכנס בתקציב. בפועל ריצה ארכה כ־15 דקות על מעבד, ולכן כל האימונים רצו מקומית.

\subsection{שלב S4: הניסוי המרכזי תחת O}\label{sec:s4}
\subsubsection{הריצות}
שש ריצות (2 עולמות $\times$ 3 זרעים) ומודל N אחד, בסדר שמזווג עולמות לפי זרע, כדי שעצירה באמצע תשאיר זוגות שלמים. אקטיבציות נשמרו בכל 12 נקודות הבדיקה עבור לוח 1, ובצעדים 0 ו־5000 עבור כל הלוחות.

\subsubsection{דגימה ולוחות}
מתוך מסלולי המבחן נקבעו מראש \textbf{שמונה לוחות זרים} של 512 מסלולים. מכל מסלול נלקחת היסטוריה אחת, בעמדה 32, כך שאין שתי נקודות מאותו מסלול. הלוחות משותפים לכל השכבות, נקודות הבדיקה והמרחבים. נקודת המבט הראשית היא \code{resid_post_3}, ה־\en{residual stream} אחרי הבלוק האחרון ולפני ה־\en{LN} הסופי, במטריקה אוקלידית גולמית, בלי \en{PCA}. נאסר להשתמש ב־\en{UMAP} או \en{t-SNE} כראיה, במרחק גאודזי מבוסס \en{k-NN} (בחירת $k$ יוצרת או מוחקת קישוריות), ובממוצעים לפי המצב האמיתי (שמכניסים את מבנה העולם מבחוץ).

\subsubsection{הומולוגיה מתמידה והמדד $S$}
על כל ענן נבנית פילטרציית \en{Vietoris--Rips} מעל $\mathbb{F}_2$, בממדים 0 ו־1. הרעיון: מגדילים בהדרגה רדיוס סביב כל נקודה, ועוקבים אחרי ״חורים״ שנוצרים ונסגרים. חור שחי לאורך טווח רדיוסים רחב הוא מבנה יציב, וחור קצר הוא רעש (איור~\ref{fig:ph}). המדד הראשי הוא אורך הפס הסופי הארוך ביותר ב־$H_1$, מנורמל בסקאלת הענן:
\begin{equation}
S=\max_{(u,v)\in D_1,\;v<\infty}\frac{v-u}{s},\qquad s=q_{95}\big(\{d(x,x')>0\}\big),
\end{equation}
ו־$S=0$ אם אין פס $H_1$ סופי.

\fig[0.9\linewidth]{ph}{הומולוגיה מתמידה (המחשה). ענן שמסודר סביב מעגל מייצר פס $H_1$ ארוך אחד ופסים קצרים רבים. הנרמול בסקאלה מאפשר להשוות עננים בגדלים שונים, למשל אקטיבציות באתחול ובסוף האימון.}

\subsubsection{הניגוד $D$ ומעטפת הרעש}
לכל חזרה $r$ ולכל לוח $i$, עם \textbf{אותו אינדקס לוח} בשני העולמות, מחושב הפרש כפול (איור~\ref{fig:dcontrast}):
\begin{equation}
D_{r,i}=\big(\Sc[,i]{cyl}{final}-\Sc[,i]{cyl}{init}\big)-\big(\Sc[,i]{sq}{final}-\Sc[,i]{sq}{init}\big).
\end{equation}
\fig[0.9\linewidth]{dcontrast}{מבנה הניגוד $D$. ההפרש בתוך כל עולם מנכה את מה שהארכיטקטורה מייצרת כבר באתחול; ההפרש בין העולמות מנכה את מה שהאימון מייצר בכל עולם. ארבעת ערכי $S$ מדווחים תמיד לצד $D$, כדי ש־$D$ לא יסתיר את מקור ההבדל.}

מעטפת הרעש מודדת כמה $S$ משתנה רק בגלל בחירת הדגימה:
\begin{equation}
u_r=\max_{c}\;q_{95}\big\{|S_{c,i}-S_{c,j}|:i<j\big\},\qquad c\in\{\text{cyl-final},\text{cyl-init},\text{sq-final},\text{sq-init}\}.
\end{equation}
הבחירה למדוד אותה מהנתונים, ולא מבקרת N, נועדה למנוע מעגליות שהופיעה באחת הטיוטות המוקדמות.

\subsubsection{הקשר: רצפה גאוסית ו־H0}
לכל ענן חושב גם $S$ של 20 עננים גאוסיים בגודל ובממד זהים, איזוטרופיים ומותאמי קווריאנס. הם אינם סף להכרעה, אלא הקשר: עבודות קודמות מראות שענן גאוסי בגודל כזה מייצר מאות פסי $H_1$ מלאכותיים, ולכן אסור לספור פסים בלי ייחוס. מודל N נבדק לזליגה: ה־\en{NLL} שלו אסור שיהיה נמוך מ־$H(W)=\ln 36$ מעבר לרווח הסמך.

\subsubsection{H2: גאומטריה של היסטוריה מול סמל אחרון}
על כל זוגות ההיסטוריות בלוח מחושב מתאם \en{Spearman} בין מרחקי האקטיבציות $d_A$ למרחקי האמונה (במטריקת \en{Hellinger}, $d_H(p,q)=\|\sqrt p-\sqrt q\|_2/\sqrt2$):
\begin{equation}
\dgeom=\rho_s\big(d_A,d_{B32}\big)-\rho_s\big(d_A,d_{B1}\big).
\end{equation}
החישוב נעשה פעם על כל הזוגות, ופעם רק על זוגות \textbf{עם אותו סמל אחרון}. בזוגות כאלה $\bel{1}$ זהה, ולכן כל התאמה לאמונה חייבת לבוא ממידע מוקדם יותר בהיסטוריה. ההכרעה היא לפי עקביות הסימן ב־3 חזרות ו־8 לוחות.

\subsubsection{בסיסי ייצוג}
ארבעה ייצוגים חלופיים עוברים \textbf{בדיוק את אותו צינור}, כל אחד במקום נקודת הבדיקה הסופית: \textbf{R0}, המודל באתחול; \textbf{R1}, ה־\en{embedding} של הסמל האחרון במודל המאומן; \textbf{R2}, ממוצע $\tfrac12[\phi(w_{32})+\phi(w_{31})]$ של \en{embedding} ספקטרלי $\phi$ של מטריצת המעברים בין סמלים $E^{\top}\mathrm{diag}(\pi)TE$; \textbf{R3}, הממוצע $\tfrac1{32}\sum_t\phi(w_t)$ על כל ההקשר, מה ש־\en{attention} כמעט אחיד מחשב. R2 ו־R3 נבנים בלי מודל. אם הם משחזרים את האות, אין צורך בסינון בייסיאני כדי להסביר אותו.

\subsubsection{ניתוח משני \en{X4}: נקודות מבט על הייצוג}
ניתוח שנרשם מראש בנפרד בדק את כלל F1 בשש נקודות מבט: V0 (הגולמית), V1 (אחרי נרמול \en{LN} בלי הפרמטרים הנלמדים), V2 (מוצא ה־\en{LN} הסופי), V3 (הרכיב שה־\en{unembedding} רואה), V4 (הרכיב שה־\en{unembedding} אינו רואה), ו־V5 (התפלגות החיזוי עצמה). המטרה היא לשאול \textbf{איפה} המבנה יושב, ולא רק אם הוא קיים.

\subsection{שלב S5: הבקרה המותאמת P}
שש ריצות נוספות, זהות לחלוטין לריצות S4 פרט למטריצת הפליטה. ההיגיון: אם המבנה ב־O נלמד מהדינמיקה ומההיסטוריה, הוא אמור לשרוד גם כשהפליטה אינה חלקה. אם הוא רק העתקה של גאומטריית הפליטה, הוא אמור להיעלם.

כבר בשלב S2 התגלה שהבקרה הזו \textbf{חד־צדדית}. תחת P, גם האמונה המדויקת $\bel{32}$ אינה מבחינה באופן יציב בין העולמות (5 מתוך 8 לוחות). לכן גם מנבא אידאלי היה נכשל בכלל F1 תחת P. כישלון תחת P לא יכול ללמד שהמבנה בא מהפליטה; רק הצלחה תחת P הייתה אינפורמטיבית. הקריאה החד־צדדית הזו נכתבה ברישום המוקדם \textbf{לפני} שהורצו ריצות P.

כדי שהשלב לא יישאר בלי שאלה אינפורמטיבית, הוצע כלל נוסף, \textbf{H3$'$}, ואושר על ידי הבעלים בנקודה H3 (החלטה \en{D-037}), לפני שנראו תוצאות P: תחת P, בצעד 5000, $\rho_s(d_A,d_{B32})$ צריך לעלות גם על ערכו באתחול וגם על $\max(\mathrm{R2},\mathrm{R3})$, בכל 3 החזרות ובלפחות 7 מתוך 8 לוחות. H3$'$ אינו שואל על חור; הוא שואל אם הגאומטריה של המודל מתיישרת עם האמונה הבייסיאנית.

\subsection{שלב S6: דינמיקת האימון}
H4 שואלת על \textbf{סדר}: האם המבנה מופיע לפני שהמודל משתמש בהיסטוריה, או אחרי. הכלל נקבע (\en{D-034}) לפני כל תוצאה של S4. לכל ריצה, על לוח 1 וב־12 נקודות הבדיקה:
\begin{equation}
t_h=\min\{t:\mathrm{regret}_t<\GC\},\qquad t_s=\min\{t:S_t-S_0>u_r\}.
\end{equation}
$t_h$ היא נקודת הבדיקה הראשונה שבה המודל כבר עדיף על כל מנבא שמשתמש רק בסמל האחרון. ריצת גליל ״עקבית״ אם $t_s\ge t_h$. H4 נתמכת אם שלוש ריצות הגליל תחת O עקביות, ונדחית אם אחת מהן מפרה את הסדר. ריצות הריבוע ו־P מדווחות תיאורית.

שני ניתוחים תיאוריים נוספים: \textbf{\en{null} של סדר נקודות הבדיקה} (1000 ערבובים, לבדוק אם השינוי ב־$S$ מרוכז בתחילת האימון יותר ממה שסדר אקראי היה נותן), ו\textbf{סריקת שכבות}, אותו ניגוד F1 בכל \en{hook} ובשרשור של כולם.

\subsection{שלב S7: בקרות, \en{probes}, קואורדינטות מעגליות ורגישות}
כל הפרטים התפעוליים של S7 הוקפאו (\en{D-047}) לפני כל ריצה או ניתוח בשלב.

\subsubsection{בקרות U ו־N}
שני מודלי U (מצב גלוי, זרע 1729) ושני מודלי N נוספים. ב־U הסמל האחרון חושף את המצב, ולכן $\GC=0$ ושער הלמידה אינו מוגדר. במקומו נבדק $\mathrm{captured}_U$ מול התועלת על פני התפלגות הסמלים השולית, $G_U=\KL(q_C\,\|\,\pi E)$. ענן האמונות של U הוא קודקודי הסימפלקס, ולכן U אינה בקרה חיובית ל־PH.

\subsubsection{\en{Probes} לינאריים}
רגרסיית \en{ridge} מהאקטיבציות אל $\bel{32}$. ההתאמה נעשית על לוחות פיתוח 1–6, הפרמטר $\lambda$ נבחר על לוחות פיתוח 7–8, והציון נמדד פעם אחת בלבד על 8 לוחות המבחן. המטרה מושווית לארבעה מתחרים: אקטיבציות באתחול; קידוד \en{one-hot} של הסמל האחרון (החסם העליון לכל מפענח שמשתמש רק בו); $F_1$ המדויק (כמה מהאמונה ניתן לשחזר ממה שמשימת החיזוי דורשת); ותוויות מעורבבות (בקרה שהפרוטוקול אינו מייצר התאמה יש מאין).

\subsubsection{קואורדינטות מעגליות}
ממחלקת ה־$H_1$ הבולטת בענן האקטיבציות מחושב קוציקל (מעל $\mathbb{F}_{47}$, באמצע חיי הפס), וממנו קואורדינטה מעגלית $\varphi$ לכל היסטוריה. ההתאמה לזווית הסמויה $\theta=2\pi x(z_{32})/12$ נמדדת כ־
\begin{equation}
A=\max\Big(\big|\overline{e^{\,i(\varphi-\theta)}}\big|,\;\big|\overline{e^{\,i(\varphi+\theta)}}\big|\Big)\in[0,1].
\end{equation}
המדד אינו רגיש לסיבוב ולשיקוף של המעגל. ההשוואות שנקבעו מראש: אתחול, ציר $y$ הלא־מחזורי, ריצות הריבוע, ותקרה מהאמונה המדויקת. זה הניתוח היחיד שמשתמש במצב האמיתי, ולכן הוא מסומן ״הערכה מיוחסת״ ותיאורי בלבד.

\subsubsection{רגישות}
שינוי של גורם אחד בכל פעם: מטריקת \en{chord} (מרחק בין וקטורים מנורמלים), \en{PCA} ל־32 ממדים שהותאם על נתוני הפיתוח, לוחות של 256 היסטוריות במקום 512, ו־$\sqrt{\mathrm{JSD}}$ במקום \en{Hellinger} לאמונות. וריאנטים אלה בודקים עמידות של הכרעה, ולעולם אינם מחליפים אותה.

\subsection{שלב S8: שחזור נקי}
שכפול טרי של המאגר ב־\en{commit} הנוכחי, עם מאגר תוצרים \textbf{ריק}, והרצה מחדש של S0 עד S3, כולל אימון מלא של ריצה אחת. כל תוצר שנוצר מושווה למאגר המקורי לפי מזהה ולפי \en{checksum} של המטען. בנוסף נבדקים שלמות המאגר כולו, שחזור הדוחות \en{E0--E3} מרשומות המדדים בלבד, וכל חבילות הבדיקות. השער \en{G-H} דורש שכל החלקים יעברו. רשומות שמכילות זמן ריצה נמדד מושוות בלי שורות הזמן (\en{D-048}).

\subsection{מתודולוגיה של מחקרי ההמשך}
אחרי \en{E4} נשארו שתי שאלות פתוחות: מקור המבנה (היסטוריה מול סמל אחרון) ושימוש סיבתי. מחקרי ההמשך נבנו לפי תבנית קבועה. ראשית נכתב פרוטוקול ב\textbf{שני מסלולי תכנון עצמאיים}, שמוזגו לטיוטה אחת. אחר כך הפרוטוקול \textbf{הוקפא} עם \en{hash}, לפני כל חישוב. ההרצה עוברת סדרה של \textbf{שערי היתכנות}, והראשון שנכשל עוצר את הכול, בלי ניסיון חוזר ובלי התאמה מחדש. לבסוף, לכל תוצאה נכתב ״רישיון״: משפט שמגדיר בדיוק מה מותר לטעון ממנה ומה אסור. תוצאות רבות מסתיימות בשער היתכנות, וזה מכוון: עדיף לעצור עם ״לא מכריע״ מאשר לפרש מדידה שאינה תקפה.

% =========================================================================================
\section{תוצאות הניסויים ומשמעותן}\label{sec:results}

\subsection{\en{E0}: הצינור מזהה את הטופולוגיה הידועה}
\fig[0.85\linewidth]{e0}{$S$ של העולמות עצמם. בשתי המטריקות לגליל יש פס $H_1$ ארוך, ולריבוע אין. הקווים המקווקווים הם ספי השער \en{G-A} שנקבעו מראש (החלטה \en{D-016}).}

\begin{table}[H]\small\centering
\begin{tabular}{@{}lcccc@{}}
\toprule
\textbf{עולם} & \textbf{מחזורי גרף} & \textbf{$\beta_1$ מעל $\mathbb{F}_2$} & \textbf{$S$, מטריקה שטוחה} & \textbf{$S$, מטריקת גרף}\\
\midrule
ריבוע & 121 & 0 & 0.037 & 0.067\\
גליל & 133 & 1 & 0.299 & 0.231\\
\bottomrule
\end{tabular}
\caption{אמת המידה: מחזורי גרף מול מספר בטי ו־$S$.}\label{tab:e0}
\end{table}

בכיול הסינתטי (512 נקודות ב־$\mathbb{R}^{128}$) ה־$S$ הממוצע של מעגל היה 0.460, וה־$S$ המרבי של דיסק 0.045; כל 20 ההגרלות סווגו נכון. השער \en{G-A} עבר.

\begin{meaningbox}[משמעות]
זו התוצאה שמתירה לפרש את כל השאר. היא מראה שני דברים. ראשית, המדד $S$ מבחין בין עולם עם חור לעולם בלי חור, עם מרווח גדול מהספים. שנית, הוא מבחין \textbf{למרות} שלשני הגרפים יש יותר מ־120 מחזורים: מה שנמדד הוא חור במשטח, ולא ספירת מחזורים בגרף. אילו השער היה נכשל, כל ניתוח של אקטיבציות היה כשל מדידה.
\end{meaningbox}

\subsection{\en{E1}: ההבדל קיים בעיקרון, עוד לפני כל מודל}\label{sec:e1}
\fig{e1}{חציון $S$ בשמונה לוחות, במרחבים שאינם נלמדים (חזרה 1729, $\sigma=1.25$, מרחקי \en{Hellinger}). מימין תחת O, משמאל תחת P. הרצפה הגאוסית לכל היותר 0.10.}

תחת O כל מרחבי הייחוס מבחינים בין העולמות ב־8 מתוך 8 לוחות, פי 3 עד 8 מהרצפה הגאוסית (איור~\ref{fig:e1}). האמונה המדויקת $\bel{32}$ מראה את הניגוד החזק ביותר (0.671 מול 0.321), וגם הוא חזק בהרבה מזה של אמונת הסמל האחרון (0.324 מול 0.165). תחת P התמונה מתהפכת: האמונה המדויקת כמעט אינה מבחינה (0.323 מול 0.310, 5 מתוך 8 לוחות), ובמרחבי העתיד הכיוון אף מתהפך (1 ו־2 מתוך 8). אמונת הסמל האחרון ושדה הפליטה זהים בשני התנאים, כי P רק מחליפה תוויות של מצבים, ולכן ענן הנקודות שלהם זהה עד כדי שינוי שמות.

\begin{meaningbox}
\begin{itemize}[leftmargin=*]
\item \textbf{ש1 ניתנת להכרעה תחת O.} יש חתימה טופולוגית צפויה, ומנבא אידאלי היה מציג אותה. אילו גם תחת O המרחבים לא היו מבחינים, הפרויקט היה נעצר כאן לפי מסלול הקיצוץ שנקבע מראש.
\item \textbf{ההיסטוריה מחזקת את החתימה.} הניגוד ב־$\bel{32}$ גדול מזה של $\bel{1}$, ולכן ל־H2 יש מה למדוד.
\item \textbf{H3 חד־צדדית.} כיוון שתחת P גם האמונה המדויקת אינה מבחינה, כישלון של המודל תחת P לא יוכל להיות ראיה נגדו. זו הייתה המסקנה המתודולוגית החשובה ביותר של השלב, והיא נרשמה לפני ריצות P.
\end{itemize}
\end{meaningbox}

\subsection{\en{E2}: הניסוי המרכזי תחת O}\label{sec:e2}
\subsubsection{המודלים למדו}
כל שש הריצות עברו את שער הלמידה בהפרש גדול: $\mathrm{captured}$ של 0.985–0.986 בגליל ו־0.992 בריבוע, כלומר \en{regret} של 0.002–0.004 \en{nats} לעומת $\GC$ של 0.24–0.26. זמן ריצה היה 13–16 דקות על מעבד.

\subsubsection{F1: נתמך}
\fig{f1panels}{הניגוד $D$ לפי לוח. כל נקודה היא לוח אחד (8 לכל חזרה). הקו השחור הוא החציון והקו האדום המקווקו הוא מעטפת הרעש $u_r$. תחת O (מימין) כל 24 הלוחות חיוביים והחציון גבוה מהמעטפת בכל חזרה. תחת P (משמאל) החציון נמוך מהמעטפת בכל חזרה (סעיף~\ref{sec:e3p}).}

\begin{table}[H]\small\centering
\begin{tabular}{@{}cccccccc@{}}
\toprule
\textbf{חזרה} & \textbf{$\Sc{cyl}{init}$} & \textbf{$\Sc{cyl}{final}$} & \textbf{$\Sc{sq}{init}$} & \textbf{$\Sc{sq}{final}$} & \textbf{לוחות $D>0$} & \textbf{חציון $D$} & \textbf{$u_r$}\\
\midrule
1729 & 0.058 & 0.307 & 0.062 & 0.129 & 8/8 & 0.158 & 0.097\\
2718 & 0.057 & 0.321 & 0.058 & 0.137 & 8/8 & 0.176 & 0.084\\
3141 & 0.058 & 0.304 & 0.059 & 0.126 & 8/8 & 0.173 & 0.084\\
\bottomrule
\end{tabular}
\caption{כלל F1 תחת O: חציון $S$ בשמונה לוחות, \code{resid_post_3}, צעד 0 מול צעד 5000.}\label{tab:f1}
\end{table}

ארבעת ערכי $S$ מראים מאיפה בא הניגוד (טבלה~\ref{tab:f1}): באתחול שני העולמות זהים (כ־0.06), ובסוף הגליל מגיע לכ־0.31 והריבוע לכ־0.13. גם הריבוע עולה, אבל הרבה פחות. ה־$S$ של הגליל בלוח 1 (0.23–0.27) גבוה פי שניים ויותר מהאחוזון ה־95 של הרצפה הגאוסית מותאמת הקווריאנס (0.10–0.12).

\subsubsection{בסיסי הייצוג: מגבלת R1}
\fig[0.9\linewidth]{baselines}{כמה מהניגוד משחזר כל בסיס: היחס בין חציון $D$ של הבסיס לחציון $D$ של המודל. ה־\en{embedding} של הסמל האחרון (R1) משחזר 0.47–0.89 מהניגוד; בסיסי הערבוב R2 ו־R3 לכל היותר 0.20.}

בסיסי הערבוב הפשוטים אינם מסבירים את האות, אבל ה־\en{embedding} של הסמל האחרון לבדו משחזר חלק ניכר ממנו (איור~\ref{fig:baselines}). כלומר, תחת O, חלק גדול מהמבנה כבר קיים ברמת ה־\en{embedding} של סמל בודד, עוד לפני שהמודל מערבב מידע בין עמדות.

\subsubsection{H2: עקבי בסימן}
\begin{table}[H]\small\centering
\begin{tabular}{@{}lcccc@{}}
\toprule
\textbf{עולם} & \textbf{חציון $\dgeom$, כל הזוגות} & \textbf{לוחות $>0$} & \textbf{חציון $\dgeom$, אותו סמל אחרון} & \textbf{לוחות $>0$}\\
\midrule
גליל (3 חזרות) & 0.124–0.127 & 24/24 & 0.935–0.939 & 24/24\\
ריבוע (3 חזרות) & 0.087–0.090 & 24/24 & 0.977–0.979 & 24/24\\
\bottomrule
\end{tabular}
\caption{H2 תחת O.}\label{tab:h2}
\end{table}
על כל הזוגות, מרחקי האקטיבציות מתואמים עם אמונת 32 הסמלים יותר מאשר עם אמונת הסמל האחרון, ב־48 מתוך 48 לוחות (טבלה~\ref{tab:h2}). בתוך זוגות עם אותו סמל אחרון, שבהם R1 קבוע ואינו יכול להסביר דבר, המתאם עם האמונה המלאה גבוה מאוד (0.94–0.98) ועולה על R2 ו־R3 בכל לוח (\en{X4-d}).

\subsubsection{בקרות ו־\en{X4}}
מודל N: \en{NLL} של 3.5835, שווה ל־$H(W)=\ln 36$ בתוך רווח סמך של 0.0002. אין זליגה, והשער \en{G-F} עבר. בניתוח \en{X4}, כלל F1 נתמך בכל שש נקודות המבט, כולל V4, הרכיב שה־\en{unembedding} אינו רואה. הבקרה החיובית של \en{X4-b} (התאמה של התפלגות החיזוי ל־$F_1$ המדויק) עמדה על 0.999 בכל ריצה.

\begin{meaningbox}
\begin{itemize}[leftmargin=*]
\item \textbf{ש1: כן, תחת O.} אימון לחיזוי הסמל הבא יוצר בייצוג אות $H_1$ שמבחין בין גליל לריבוע, עקבי בכל לוח ובכל חזרה, מעל הרעש ומעל הרצפה הגאוסית.
\item \textbf{ש2, חלק (ב): יש שימוש בהיסטוריה, בגאומטריה.} H2 מראה שהמרחקים בין ייצוגים עוקבים אחרי האמונה המלאה, גם כשהסמל האחרון זהה.
\item \textbf{ש2, חלק (א): לא ניתן לשלול את גאומטריית הסמל.} R1 נושא עד 89\% מהניגוד. F1 לבדו אינו מראה ששילוב ההיסטוריה הוא שיוצר את החור. לכן ההכרעה במקור עברה ל־P.
\item \textbf{מיקום:} המבנה קיים גם ברכיב הנראה לפלט וגם ברכיב הנסתר ממנו (\en{X4-c}: ״\en{both}״). זה מראה שהוא קיים שם, לא שהוא נחוץ לחיזוי.
\end{itemize}
\end{meaningbox}

\subsection{\en{E3}: הבקרה המותאמת P}\label{sec:e3p}
כל שש ריצות P עברו את שער הלמידה ($\mathrm{captured}$ 0.874–0.914). ובכל זאת כלל F1 \textbf{נדחה} תחת P: חציון $D$ של 0.021–0.034, מתחת למעטפת הרעש (0.066–0.089) בכל חזרה, ובשתי חזרות גם פחות מ־7 לוחות חיוביים (איור~\ref{fig:f1panels}, משמאל). ה־\en{embedding} של הסמל האחרון, שנשא את רוב הניגוד תחת O, קורס תחת P (חציון $D$ מ־$0.140$ ל־$-0.035$ בחזרה 1729).

לעומת זאת, \textbf{H3$'$ נתמך}: תחת P, $\rho_s(d_A,d_{B32})$ עולה מ־0.06–0.11 באתחול ל־0.70 בסוף האימון, ועולה על הטוב מבין R2 ו־R3 (0.39) ב־48 מתוך 48 לוחות. גם בתוך זוגות עם אותו סמל אחרון (\en{X4-d}) המודל עולה על בסיסי הערבוב, עם מתאם של כ־0.82.

\begin{table}[H]\small\centering
\begin{tabularx}{\linewidth}{@{}Y cc@{}}
\toprule
\textbf{פריט} & \textbf{תחת O} & \textbf{תחת P}\\
\midrule
כלל F1, בכל אחת מנקודות המבט \en{V0--V5} & \good{נתמך} & \bad{נדחה}\\
H3$'$ (התיישרות עם האמונה), V0 ו־V1 & \good{נתמך} (השוואה) & \good{נתמך} (הכרעה)\\
\en{X4-d} (אותו סמל אחרון, מול R2/R3) & \good{נתמך} & \good{נתמך}\\
\en{X4-c} (איפה המבנה יושב) & בשני הרכיבים & באף אחד\\
חציון $\dgeom$, כל הזוגות & 0.088–0.125 & 0.041–0.072\\
\bottomrule
\end{tabularx}
\caption{O מול P.}\label{tab:op}
\end{table}

\begin{meaningbox}
\begin{itemize}[leftmargin=*]
\item \textbf{הדחייה של H3 אינה ראיה לכך שהמבנה בא מהפליטה.} כפי שנרשם מראש, גם מנבא מושלם היה נכשל כאן, כי האמונה המדויקת עצמה אינה נושאת את החור תחת P. התוצאה פשוט אינה עונה על השאלה.
\item \textbf{התיישרות עם אמונה $\neq$ ניגוד $H_1$.} H3$'$ מראה שהמודל לומד גאומטריה שעוקבת אחרי האמונה הבייסיאנית גם כשאין חור לזהות. ניגוד $H_1$ צפוי רק כאשר גאומטריית האמונה עצמה נושאת את החור של העולם. זו אחת התובנות המרכזיות של הפרויקט.
\item \textbf{ש2 נשארת פתוחה.} הקריסה של R1 תחת P מראה שהניגוד של R1 תחת O אינו גאומטריית פליטה סטטית, כי P שומרת על ענן הפליטה. אבל היא אינה מבחינה בין שילוב היסטוריה לבין סטטיסטיקת מעברים שנלמדה ברמת הסמל.
\end{itemize}
\end{meaningbox}

\subsection{\en{E3}: דינמיקת האימון (H4)}
\fig{dynamics}{$S$ ו־\en{regret} לאורך האימון, לוח 1. כל קו הוא ריצה. בשני התנאים, $S$ קופץ כבר בצעד 50 בשני העולמות, וה־\en{regret} יורד מתחת ל־$\GC$ רק בצעד 100. תחת O (מימין) הפער בין הגליל לריבוע נפתח בעיקר בין צעד 100 לצעד 200.}

בכל שלוש ריצות הגליל תחת O, $t_h=100$ ו־$t_s=50$ (איור~\ref{fig:dynamics}). לפי הכלל, ריצה שבה $S$ עולה לפני השימוש בהיסטוריה מפרה את H4, ולכן \textbf{H4 נדחתה}. דפוס זהה מופיע בכל ריצות הריבוע תחת O, ודפוס דומה תחת P (שם רק ריצה אחת, ריבוע 2718, עקבית עם הכלל). ה־\en{null} של סדר נקודות הבדיקה אינו מכריע ($p$ בין 0.07 ל־0.18 בריצות הגליל).

\fig[0.9\linewidth]{layers}{סריקת שכבות: חציון $D$ לכל \en{hook}. תחת O הניגוד קיים כבר ב־\en{embedding} (כ־0.11) וגדל לאורך השכבות עד ל־\en{LN} הסופי (כ־0.28). תחת P הוא קרוב לאפס בכל השכבות. תיאורי בלבד.}

\begin{meaningbox}
\begin{itemize}[leftmargin=*]
\item \textbf{העלייה המוקדמת אינה ייחודית לגליל.} בצעד 50 גם הריבוע עולה (0.17–0.23, לעומת 0.24–0.25 בגליל). זו עלייה כללית של מבנה בענן, ולא הופעה של החור.
\item \textbf{הניגוד עצמו מופיע מאוחר יותר.} ההפרש בין הגליל לריבוע בלוח 1 הוא 0.02–0.07 בצעד 50 ו־0.18–0.23 בצעד 200, כלומר אחרי שהשימוש בהיסטוריה מתחיל. אבל זו תצפית חקרנית: הכלל שנרשם מראש התייחס ל־$S$ הגולמי של הגליל, ולא לניגוד. היא אינה משנה את דחיית H4.
\item \textbf{הניגוד קיים כבר ב־\en{embedding}} (איור~\ref{fig:layers}). זה מחזק את מגבלת R1: תחת O חלק מהמבנה נוצר ברמת הסמל הבודד, עוד לפני כל ערבוב בין עמדות.
\end{itemize}
\end{meaningbox}

\subsection{\en{E4}: בקרות U ו־N}
שני מודלי N נוספים: ה־\en{NLL} שווה ל־$\ln 36$ בתוך $2\cdot10^{-5}$, וה־$S$ של האקטיבציות נשאר ברמת רעש (0.002–0.077). מודלי U למדו את הדינמיקה כמעט במלואה ($\mathrm{captured}_U=0.999$, \en{regret} של 0.005 \en{nats} מול תועלת של 3.4 \en{nats}). ב־U הגליל הגיע ל־$S$ של 0.215 והריבוע ל־0.147, עם $D$ חיובי ב־8 מתוך 8 לוחות.

\begin{meaningbox}[משמעות]
N מאשרת שאין זליגה ושהארכיטקטורה לבדה אינה מייצרת את האות. U מראה משהו חשוב לשאלת המקור: כשהמצב גלוי, ואין שום צורך בהסקה רב־שלבית, עדיין נוצר בייצוג מעגל שמיושר עם ציר $x$ (סעיף~\ref{sec:circ}). כלומר, \textbf{סטטיסטיקת מעברים ברמת הסמל יכולה להספיק} כדי ליצור מעגל כזה. זו ראיה שהמסלול החלופי קיים, לא שהוא המסלול שפעל תחת O.
\end{meaningbox}

\subsection{\en{E4}: \en{probes} לינאריים}
\fig{probes}{מה אפשר לפענח לינארית ($R^2$ על המבחן, חציון על 6 ריצות לכל תנאי). שתי השורות העליונות כמעט זהות: האקטיבציות אינן מפענחות את האמונה טוב יותר מהתפלגות הסמל הבא המדויקת.}

האימון מעלה את יכולת הפענוח של האמונה הרבה מעל האתחול (0.83 מול 0.48 תחת O) ומעל כל מפענח של הסמל האחרון (0.42). בקרת התוויות המעורבבות נמצאת באפס, כך שהפרוטוקול אינו מייצר התאמה. אבל הדפוס התיאורי שנרשם מראש, שהאקטיבציות יעלו על כל שלושת המתחרים, הופיע רק בריצה אחת מתוך שש בכל תנאי, כי הן אינן עולות על $F_1$ (איור~\ref{fig:probes}). האקטיבציות מפענחות את אמונת הסמל האחרון כמעט בשלמות (0.97).

\begin{meaningbox}
לגבי ש2(ג): \textbf{במובן הלינארי, הייצוג נושא בערך את מה שמשימת החיזוי דורשת, ולא יותר.} אין כאן ראיה לתוכן אמונה מעבר לצורכי הסמל הבא. זה אינו סותר את H2 ואת H3$'$, שעוסקים בגאומטריה של מרחקים ולא בפענוח לינארי, אבל זה מסייג אותם: התאמה לגאומטריית האמונה מתיישבת גם עם קידוד של מה שדרוש לחיזוי בלבד. הפענוח לינארי בלבד, ולכן רכיבים לא־לינאריים אינם נשללים.
\end{meaningbox}
\begin{caveat}
תחת O, ה־$\lambda$ שנבחר היה הערך הקטן ביותר ברשת ($10^{-3}$) בכל הריצות, כך שהרשת לא תחמה את האופטימום. סביר שההשפעה קטנה, אבל $R^2$ תחת O עשוי להיות מוערך מעט בחסר. הרשת הוקפאה ב־\en{D-047} ולא שונתה בדיעבד.
\end{caveat}

\subsection{\en{E4}: קואורדינטות מעגליות}\label{sec:circ}
\fig{circular}{התאמת הקואורדינטה המעגלית לזווית הסמויה של ציר $x$ (חציון ההתאמה $A$). מודלי הגליל המאומנים תחת O מגיעים ל־0.89, כ־97\% מהתקרה של האמונה המדויקת (0.915).}

תחת O, האימון הופך את מחלקת ה־$H_1$ הבולטת של הגליל לקואורדינטה על ציר המחזור של העולם: התאמה של 0.884–0.892 בשלוש הריצות, לעומת 0.06–0.12 באתחול, 0.10 מול ציר $y$ הלא־מחזורי, ו־0.21–0.30 בריצות הריבוע. תחת P ההתאמה נשארת ברצפה (0.05).

\begin{meaningbox}
זו הראיה הישירה ביותר שהאות שזוהה ב־F1 הוא \textbf{המעגל של הגליל}, ולא לולאה אחרת כלשהי בענן. היא מזהה \textbf{איזו} לולאה F1 מודד. היא אינה ראיה לשימוש סיבתי, ואינה מבחינה בין שילוב היסטוריה למבנה של הסמל האחרון: לסטטיסטיקה הזו לא היה קו בסיס של סמל אחרון (נבדק מאוחר יותר, סעיף~\ref{sec:origin}), ותוצאת U מראה שמעגל כזה יכול לנבוע גם מסטטיסטיקה ברמת הסמל.
\end{meaningbox}
\begin{caveat}[הסתייגויות]
(1) ה־\en{null} של ערבוב הזוויות חלש: גם באתחול 15 מתוך 24 לוחות עוברים את האחוזון ה־95 שלו. ההשוואות האינפורמטיביות הן אלה שנקבעו מראש. (2) הקואורדינטות מחושבות בסקאלת האמצע של הפס; אם השלד בסקאלה זו אינו קשיר, רכיבים מחוץ לרכיב הענק מקבלים היסט שרירותי. אבחון גאומטרי מאוחר (משימה 0023) מצא לכל היותר נקודה אחת בכל לוח מחוץ לרכיב הענק. הערכים שדווחו עומדים.
\end{caveat}

\subsection{\en{E4}: רגישות}
\begin{table}[H]\small\centering
\begin{tabularx}{\linewidth}{@{}Y cccc@{}}
\toprule
\textbf{וריאנט (גורם אחד)} & \textbf{F1 תחת O} & \textbf{F1 תחת P} & \textbf{H3$'$ תחת P} & \textbf{שיעור $\dgeom>0$ (O / P)}\\
\midrule
ראשי (אוקלידי גולמי, $M=512$) & \good{נתמך} & \bad{נדחה} & \good{נתמך} & \en{1.00 / 1.00}\\
מטריקת \en{chord} & \good{נתמך} & \bad{נדחה} & \good{נתמך} & \en{1.00 / 0.90}\\
\en{PCA-32} מותאם על \en{dev} & \good{נתמך} & \bad{נדחה} & \good{נתמך} & \en{1.00 / 1.00}\\
256 היסטוריות ללוח & \good{נתמך} & \bad{נדחה} & \good{נתמך} & \en{1.00 / 1.00}\\
$\sqrt{\mathrm{JSD}}$ לאמונות & ללא שינוי & ללא שינוי & \good{נתמך} & \en{1.00 / 0.71}\\
\bottomrule
\end{tabularx}
\caption{וריאנטי רגישות, גורם אחד בכל פעם.}\label{tab:sens}
\end{table}
\begin{meaningbox}[משמעות]
ההכרעות הראשיות אינן תלויות באף אחת מארבע הבחירות האנליטיות שנבדקו (טבלה~\ref{tab:sens}). זו עמידות מול ארבעה שינויים מסוימים, לא מול כל בחירה אפשרית: ה־\en{hook}, נקודת הבדיקה, הפילטרציה, שדה המקדמים, $\sigma$ וגודל העולם לא שונו.
\end{meaningbox}

\subsection{\en{E4}: שחזור (S8)}\label{sec:repro}
\begin{table}[H]\small
\begin{tabularx}{\linewidth}{@{}>{\raggedleft\arraybackslash}p{0.3\linewidth} Y c@{}}
\toprule
\textbf{בדיקה} & \textbf{תוצאה} & \textbf{מצב}\\
\midrule
תוצרים משותפים לשחזור ולמאגר המקורי & 307/307 זהים ביט־לביט, כולל 12 נקודות בדיקה ו־12 מערכי אקטיבציות של הריצה שאומנה מחדש & \good{עבר}\\
הכרעות שערים & 6/6 זהות & \good{עבר}\\
תוצר שקיים רק בשחזור & רשומת זמן אחת (\code{run_timing}), שהקוד שכותב אותה נוסף אחרי הריצה המקורית & \bad{נכשל}\\
שלמות המאגר & 785/785 תוצרים עם הורים ו־\en{checksums} תקינים & \good{עבר}\\
שחזור הדוחות \en{E0--E3} & 9/9 קבצים זהים בייט־לבייט & \good{עבר}\\
בדיקות & 511 במודולים, 3 איטיות, 4 אינטגרציה & \good{עבר}\\
\bottomrule
\end{tabularx}
\caption{השער \en{G-H}.}\label{tab:repro}
\end{table}
\begin{meaningbox}[משמעות]
האימון על מעבד דטרמיניסטי ביט־לביט בסביבה הזו, וכל המספרים בדוח ניתנים לשחזור מהרשומות. השער \en{G-H} בכל זאת \textbf{נכשל}, וההכרעה נשארה ״נכשל״ (\en{D-049}): הבדיקה שהוגדרה מראש דרשה התאמה מלאה של המאגר, והיא לא התקיימה. זה פער תיעודי שנוצר מהתפתחות הקוד, ולא הבדל מספרי, אבל אסור לקרוא אותו כ״עבר״. היקף השחזור מוגבל: אותה מכונה, אותה סביבה, ורק אחת מ־17 ריצות האימון אומנה מחדש.
\end{meaningbox}

% =========================================================================================
\section{מחקרי ההמשך אחרי \en{E4}}\label{sec:followups}
אחרי \en{E4} נותרו פתוחות שאלת המקור (היסטוריה או סמל אחרון) ושאלת השימוש. תשעה מחקרי המשך ניסו להתקדם בהן (איור~\ref{fig:followups}, טבלה~\ref{tab:followups}). אף אחד מהם אינו משנה הכרעה של \en{E0--E4}.

\fig[0.95\linewidth]{followups}{מפת מחקרי ההמשך, מקובצים לפי השאלה שכל אחד ניסה לקדם.}

\begin{table}[H]\small
\begin{tabularx}{\linewidth}{@{}Y l l@{}}
\toprule
\textbf{מחקר} & \textbf{סימון} & \textbf{תוצאה}\\
\midrule
\en{X4} שלב B: הפרדה באופק העתידי (\en{D-051--052}) & רשום מראש & \warn{לא מכריע}\\
\en{X4} שלב B2: כיסוי (\en{D-053--055}) & רשום מראש & \bad{\en{NO-GO} מוגבל}\\
ביקורת קו־הבסיס המעגלי (\en{D-056--057}) & רשום מראש, תיאורי & \neu{מעבר לסמל ולערבוב}\\
\en{X4-F}: שימוש סיבתי בקואורדינטה (\en{D-058--063}) & רשום מראש & \warn{לא מכריע}\\
פרופיל אורקל של היסטוריה מעגלית (0014, \en{D-074}) & חקרני & \neu{התקבל, מסויג}\\
מחזור מקור־ההיסטוריה K8 & רשום מראש & \bad{\en{FAIL} במרווח}\\
\en{X5 G1}: זמן למידה של זוג קרנלים (\en{D-072--073}, \en{D-081}) & רשום מראש & \bad{H5 נדחתה}\\
K4: היתכנות היסטוריה מעגלית (0016–0018, \en{D-075--078}) & רשום מראש & \warn{לא מכריע}\\
\en{D-1}: גליל מול טבעת מביוס, אורקל בלבד (0019–0020) & רשום מראש & \good{\en{PASS}}\\
\bottomrule
\end{tabularx}
\caption{מחקרי ההמשך.}\label{tab:followups}
\end{table}

\subsection{האם המבנה נחוץ לחיזוי? (\en{X4 B, B2})}
\textbf{שיטה.} השאלה כאן היא אם יש זוגות היסטוריות שנראים זהים ב־$F_1$ אבל שונים באופק 2. אם יש כאלה, מנבא אידאלי צריך לייצג יותר מהסמל הבא, ואפשר לבדוק אם המודל עושה זאת. \textbf{תוצאה.} במרחב הסובב, דרגות המשיק ($Q_1=35$, $Q_2=137$) מאפשרות הפרדה כזו בעיקרון. אבל בשלב B ארבע בריכות הדגימה לא כוסו מספיק (49–66 עוגנים כשרים מול 128 נדרשים), ולכן התוצאה לא מכריעה. בשלב B2, עם ארבע בריכות חדשות (141–159 עוגנים כל אחת), נמצאו \textbf{אפס} עדים.
\begin{meaningbox}[משמעות]
בספים וברזולוציית הדגימה שנקבעו, אין פער מעשי בין אופק 1 לאופק 2. זה אינו מוכיח ש־$F_1$ מספיק תמיד, אבל זה שולל את ההנחה שעליה נשען שלב C, ולכן הוא לא הורץ. לא נבדקו מודל, אקטיבציות או מפענח.
\end{meaningbox}

\subsection{האם המבנה בא מההיסטוריה? (ביקורת מעגלית, 0014, K8, K4)}\label{sec:origin}
\textbf{ביקורת קו־הבסיס המעגלי} סגרה, באופן תיאורי, את החסר שצוין ב־\en{E4}: ב־24 מתוך 24 לוחות, ההתאמה המעגלית של המודל (0.884–0.892) עולה על חסם אופטימי מהסמל האחרון (0.817) ועל קו הבסיס הטוב ביותר הזמין (0.805). הביקורת רצה על לוחות שכבר נחשפו, ולכן היא אינה אישוש בלתי תלוי ואינה ראיה למקור בהיסטוריה ארוכה.

\textbf{פרופיל האורקל (משימה 0014)} שאל כמה מידע על הזווית המעגלית מוסיפה היסטוריה ארוכה יותר, ומדד את ההפרש בין הפוסטריור מ־32 תצפיות לפוסטריור מ־$K$ התצפיות האחרונות (איור~\ref{fig:oracle}):
\begin{equation}
B_K=\mathbb{E}\,\big\|m_{32}-m_K\big\|,\qquad K\in\{1,2,4,8,16,32\}.
\end{equation}

\fig[0.85\linewidth]{oracle}{כמה מידע מוסיפה ההיסטוריה מעבר ל־$K$ התצפיות האחרונות (פרופיל האורקל, נתוני כיול). ההפרש יורד מהר: כבר ב־$K=8$ הוא מתחת לערך הייחוס 0.01.}

המשמעות: הפוסטריור של האורקל מתכנס מהר יחסית ל־32 תצפיות, ולכן יש ״מרווח״ מעשי לניסוי מקור־היסטוריה רק בהשוואה ל־$K$ קטן (עד 4). זה אינו אופק זיכרון של המודל או של העולם, ול־$K=8$ עדיין יש מידע שאינו אפס. \textbf{מחזור K8} ניסה בכל זאת להשוות ל־$K=8$, ושער המרווח שהוקפא נכשל: $B_8=0.00925$, עם קצה עליון חד־צדדי של 95\% ב־0.00949, מתחת ל־0.01. לאפקט של 0.01 לפחות אין מקום במערך הזה. \textbf{K4} עבר לכן ל־$K=4$ ועבר את שערי הזהות והתהליך (שגיאות מומנט עד $8\cdot10^{-16}$), אבל נעצר בשער המפה הראשון: בסקאלה שהוקפאה, גרף נקודות הציון של המודל הראשון אינו קשיר. לפי הכלל לא היה ניסיון חוזר, והניקוד לא שוחרר.

\begin{meaningbox}[משמעות לאור ש2]
שאלת המקור נשארת פתוחה, אבל היא הוגדרה בצורה חדה יותר. ראשית, המעגל שהמודל בונה טוב יותר ממה שהסמל האחרון לבדו מאפשר (תיאורית). שנית, ההיסטוריה שמעבר לכמה צעדים אחרונים כמעט אינה מוסיפה מידע על הזווית, ולכן כל ניסוי שמנסה להפריד ״היסטוריה ארוכה״ מ״קצרה״ חייב לעבוד עם $K$ קטן מאוד. שלישית, כלל המפה של K4 אינו עמיד מספיק, וזו תוצאה תפעולית, לא ראיה על האפקט עצמו.
\end{meaningbox}

\subsection{האם המודל משתמש במבנה? (\en{X4-F})}
\textbf{שיטה.} מוחקים מהאקטיבציות את הכיוון שנושא את הקואורדינטה המעגלית, ומשווים את הפגיעה בחיזוי לפגיעה של מוחקים אקראיים שגורמים ״נזק״ דומה. \textbf{תוצאה.} המוחק הממוקד ביטל את השחזור הלינארי של הקואורדינטה, אבל שחזור לא־לינארי (\en{RBF}) נשאר 0.998: המידע לא נמחק, רק הרכיב הלינארי שלו. בנוסף, אף אחד מ־10,000 מוחקים אקראיים לא נפל בחלון התאמת הנזק שהוקפא, כך שלא הייתה משפחת השוואה. נקודת הקצה הסיבתית \textbf{לא הורצה}.
\begin{meaningbox}[משמעות לאור ש3]
שאלת השימוש נשארת פתוחה לגמרי. התוצאה מלמדת דבר מתודולוגי: המידע על המעגל מקודד באופן לא־לינארי ומבוזר, ולכן מחיקה לינארית אינה מספיקה כדי לבדוק אותו. ניסוי התערבות עתידי צריך כלי אחר.
\end{meaningbox}

\subsection{הרחבות (\en{X5 G1}, \en{D-1})}
\textbf{\en{X5 G1}} בדק אם שינוי גאומטריה (קרנל ״מוליכות״ מול קרנל שטוח) מעכב את הלמידה. העיכובים המזווגים היו 250, 0 ו־0 צעדים, והטענה שנרשמה, עיכוב של לפחות 250 צעדים בכל שלושת הזוגות, נדחתה. זו אינה תוצאת אפס כללית לגאומטריה: שינוי הקרנל מעורבב עם מהירות התנועה. \textbf{\en{D-1}} בדק, באורקל בלבד ובספירה מלאה של 37,190 ההיסטוריות הישיגות, הבחנה של \textbf{אוריינטביליות} בין גליל לטבעת מביוס, שיש להן אותם מספרי בטי. כל השערים עברו: ההבחנה אינה נחוצה באופק 1 ונחוצה באופק 2. זו תוצאת קיום על המשימה, ואינה אומרת דבר על מודלים מאומנים.

% =========================================================================================
\section{דיון: מה למדנו על שאלות המחקר}\label{sec:discussion}

\subsection{ש1: האם נוצר מבנה טופולוגי? כן, תחת O}
שרשרת הראיות שלמה ועקבית. (1) הצינור מזהה את החור על אמת ידועה (\en{E0}). (2) מנבא אידאלי היה מציג את ההבדל (\en{E1}). (3) המודלים למדו (שער הלמידה). (4) כלל F1, שהוקפא מראש, מתקיים בכל לוח ובכל חזרה, מעל הרעש ומעל הרצפה הגאוסית (\en{E2}). (5) הוא עמיד לארבעה שינויי ניתוח (S7). (6) המחלקה שמזוהה היא המעגל המחזורי של העולם, קרוב לתקרה של האמונה המדויקת (S7). התשובה לש1 חיובית, בגבולות של עולם צעצוע אחד ומודל קטן אחד.

\subsection{ש2: מאין בא המבנה? התשובה חלקית}
\begin{table}[H]\small
\begin{tabularx}{\linewidth}{@{}>{\raggedleft\arraybackslash}p{0.17\linewidth} Y Y c@{}}
\toprule
\textbf{מקור אפשרי} & \textbf{ראיות בעד} & \textbf{ראיות נגד או מגבילות} & \textbf{מצב}\\
\midrule
\textbf{(א) גאומטריית התצפיות} & R1 נושא 0.47–0.89 מהניגוד; הניגוד קיים כבר ב־\en{embedding}; U מראה שסטטיסטיקה ברמת הסמל מספיקה ליצירת מעגל & P שומרת על ענן הפליטה ובכל זאת R1 קורס, כך שזו אינה גאומטריית פליטה סטטית & \warn{לא נשלל}\\
\textbf{(ב) אמונה מההיסטוריה} & H2 בכל 48 הלוחות, גם באותו סמל אחרון; H3$'$ תחת P; הקואורדינטה המעגלית עולה על חסם הסמל האחרון & H4 נדחתה; K8 נכשל במרווח; K4 נעצר; ההיסטוריה מעבר ל־8 צעדים כמעט אינה מוסיפה מידע & \warn{נתמך חלקית}\\
\textbf{(ג) רק צורכי החיזוי} & \en{probes}: האקטיבציות אינן מפענחות את האמונה טוב יותר מ־$F_1$; B2: אין פער מעשי בין אופק 1 ל־2 & \en{probes} לינאריים בלבד & \neu{מתיישב עם הנתונים}\\
\bottomrule
\end{tabularx}
\caption{מפת הראיות לשאלת המקור.}\label{tab:evidence}
\end{table}
הקריאה הזהירה של טבלה~\ref{tab:evidence}: הייצוג מתיישר עם גאומטריית האמונה הבייסיאנית, ובגליל תחת O הגאומטריה הזו נושאת את המעגל של העולם. אבל אין ראיה שהוא נושא יותר ממה שחיזוי הסמל הבא דורש, ולא הוכרע אם המבנה נבנה משילוב היסטוריה או מסטטיסטיקה ברמת הסמל. שני המסלולים אפשריים, ולשניהם יש ראיות.

\subsection{ש3: האם המודל משתמש במבנה? לא נבדק בהצלחה}
אין בפרויקט אף ראיה סיבתית. ניסיון ההתערבות היחיד (\en{X4-F}) נעצר בשער היתכנות, כי המידע מקודד לא־לינארית ולא נמצאה משפחת השוואה מותאמת. ההבחנה בין ״זיהוי״ ל״שימוש״ שנקבעה בתחילת הדרך נשארת בתוקף.

\subsection{תובנות מתודולוגיות}
\begin{itemize}[leftmargin=*]
\item \textbf{בקרה מותאמת חייבת להיבדק על האורקל לפני שמאמנים.} בלי טבלת \en{E1}, דחיית H3 הייתה נקראת כראיה שהמבנה בא מהפליטה. בפועל זו בקרה חד־צדדית.
\item \textbf{התיישרות עם אמונה וניגוד טופולוגי הן תכונות שונות.} אפשר לקבל את הראשונה בלי השנייה (P), והשנייה צפויה רק כשגאומטריית האמונה עצמה נושאת את החור.
\item \textbf{״$S$ עולה״ ו״הניגוד עולה״ הם מדדים שונים.} כלל H4 נוסח על $S$ הגולמי, והתזמון המעניין התגלה בניגוד. הניסוח של כלל קובע מה הוא יכול לגלות.
\item \textbf{שערי היתכנות חוסכים מסקנות שגויות.} רוב מחקרי ההמשך נעצרו בשער. כל עצירה כזו מנעה פרשנות של מדידה שאינה תקפה, ולימדה משהו על המשימה (התכנסות מהירה של הפוסטריור, קידוד לא־לינארי, כלל מפה שביר).
\end{itemize}

% =========================================================================================
\section{מגבלות}\label{sec:limits}
\begin{enumerate}[leftmargin=*]
\item \textbf{אין טענות סיבתיות.} כל הראיות ייצוגיות: PH של עננים, מתאמי דרגות, \en{probes} לינאריים וקואורדינטות מעגליות.
\item \textbf{מגבלת R1 עומדת בעינה.} השאלה אם המבנה בא משילוב היסטוריה או מסטטיסטיקה ברמת הסמל לא הוכרעה.
\item \textbf{P חד־צדדית.} כישלון תחת P אינו אומר דבר על למידה.
\item \textbf{מערך צעצוע אחד.} סריג $12\times12$, שתי טופולוגיות, $\sigma$ אחד, \en{Transformer} של 4 שכבות וממד 128 שאומן מאפס על רצפי \en{HMM} סינתטיים. אין כאן ראיה על מודלי שפה טבעית גדולים.
\item \textbf{סטטיסטיקה.} שלוש חזרות לכל תא. הלוחות הם חלוקות של מדגם בתוך ריצה, ולא חזרות. כלל F1 הוא עקביות בסימן ומעטפת, לא מבחן מובהקות. מודלי U הם חזרה אחת לכל עולם.
\item \textbf{בחירות PH.} \en{Vietoris--Rips} מעל $\mathbb{F}_2$, $M=512$, $S$ כפס הארוך ביותר. הרגישות כיסתה ארבעה שינויים בודדים בלבד.
\item \textbf{קואורדינטות מעגליות.} קוציקל אחד באמצע חיי הפס; המדד אינו רגיש לסיבוב ולשיקוף אבל כן לעיוות לא־לינארי של המעגל, ולכן עשוי להעריך בחסר.
\item \textbf{היקף השחזור.} אותה מכונה, מערכת הפעלה וסביבה; ריצת אימון אחת מתוך 17. לא נבדקו דטרמיניזם בין פלטפורמות ועל \en{GPU}. בדיקות האינטגרציה I03 ו־I05 לא נבנו.
\end{enumerate}

% =========================================================================================
\section{נספחים}\label{sec:appendix}
\subsection{מילון מונחים}
\begin{table}[H]\small
\begin{tabularx}{\linewidth}{@{}>{\raggedleft\arraybackslash}p{0.25\linewidth} Y@{}}
\toprule
\textbf{מונח} & \textbf{הסבר}\\
\midrule
הומולוגיה מתמידה (PH) & שיטה למדידת ״חורים״ בענן נקודות לאורך טווח של סקאלות; חור שחי לאורך טווח רחב נחשב יציב\\
$H_1$, $\beta_1$ & הומולוגיה מממד 1 (לולאות) ומספר בטי הראשון, מספר החורים החד־ממדיים הבלתי תלויים\\
$S$ & אורך פס ה־$H_1$ הסופי הארוך ביותר, מנורמל בסקאלת הענן\\
$D$, $u_r$ & הניגוד הכפול גליל־ריבוע וסוף־אתחול; ומעטפת הרעש בין לוחות\\
$\bel{32}$, $\bel{1}$ & האמונה הבייסיאנית על המצב אחרי 32 סמלים, ואחרי הסמל האחרון בלבד\\
$F_1$, $F_2$ & התפלגות הסמל הבא, וההתפלגות המשותפת של שני הסמלים הבאים\\
$\GC$, \en{regret}, \en{captured} & תועלת ההקשר; המרחק מהמנבא המדויק; והחלק מתועלת ההקשר שהמודל תפס\\
$\dgeom$ & ההפרש בין התאמת הגאומטריה לאמונה המלאה לבין התאמתה לאמונת הסמל האחרון\\
R0–R3 & בסיסי ייצוג: אתחול, \en{embedding} סמל אחרון, ערבוב עם הסמל הקודם, ממוצע הקשר\\
O, P, N, U & תנאי פליטה: חלקה, מעורבבת לפי צירים, זהה, וסמל ייחודי לכל מצב\\
\en{V0--V5} & נקודות מבט על הייצוג בניתוח \en{X4}, מהגולמי ועד התפלגות החיזוי\\
לוח (\en{panel}) & קבוצה זרה של 512 היסטוריות מבחן, אחת לכל מסלול, משותפת לכל הניתוחים\\
שער (\en{gate}) & בדיקה אוטומטית שחייבת לעבור לפני שמפרשים את תוצרי השלב\\
\bottomrule
\end{tabularx}
\caption{מילון מונחים.}\label{tab:glossary}
\end{table}

\subsection{מקורות בתוך המאגר}
\begin{itemize}[leftmargin=*]
\raggedleft
\item שאלות, השערות והחלטות סגורות: \code{docs/RESEARCH_SUMMARY.md}, \code{docs/RESEARCH_PLAN_HE.md}, \code{PREREGISTRATION.md}, \code{PREREGISTRATION_X4.md}.
\item שלבים ושערים: \code{docs/STAGES.md}. החלטות: \code{state/DECISIONS.md}.
\item תוצאות: \code{reports/E0-E4/summary.md}, \code{reports/E3/condition_P.md}, \code{reports/E4/interpretation.md}.
\item מחקרי המשך: הסעיף ״\en{After E4}״ ב־\code{README.md} והרשומות שהוא מפנה אליהן.
\item עקומות הדינמיקה (איור~\ref{fig:dynamics}) וערכי $D$ לפי לוח (איור~\ref{fig:f1panels}) נקראו ישירות מרשומות המדדים במאגר התוצרים, בקריאה בלבד (\code{extract_data.py}). שאר המספרים מועתקים מהדוחות. איורים \ref{fig:world}–\ref{fig:dcontrast} הם המחשות.
\end{itemize}

\end{document}
